\section*{Quantitative closure of the curvature step}

The only non-numerical constant needed in the run argument is a nondegenerate
interval on which the limiting density has strictly positive logarithmic
curvature.  This interval follows quantitatively from the three-window valley;
no numerical reconstruction of the limiting density is required.

Let
\[
J_L=[-q-\tfrac14,-q+\tfrac14],\qquad
J_M=[\tfrac14-q,\tfrac34-q],\qquad
J_R=[\tfrac34-q,\tfrac54-q].
\]
All three intervals have length \(1/2\).  From the variance bound
\(v_s\le 5/256\) and the component weights \(q,1-q\ge49/100\),
Chebyshev gives
\[
 \min\{\Pr(Z\in J_L),\Pr(Z\in J_R)\}\ge \frac{539}{1600},
 \qquad
 \Pr(Z\in J_M)\le\frac{5}{16}.
\]
Hence the average density on each outer interval is at least
\(539/800>67/100\), whereas the average density on the middle interval is at
most \(5/8<63/100\).  Since the limiting density \(f\) is continuous and
strictly positive, there are interior points
\[
x_L\in J_L,\qquad x_M\in J_M,\qquad x_R\in J_R
\]
such that
\[
f(x_L),f(x_R)>\frac{67}{100},
\qquad
f(x_M)<\frac{63}{100}.
\]
Put \(g=\log f\).  The point \(x_M\) lies between \(x_L,x_R\), and the value of
the affine chord through \((x_L,g(x_L))\) and \((x_R,g(x_R))\) at \(x_M\)
exceeds \(g(x_M)\) by more than
\[
\Delta:=\log\frac{67}{63}.
\]
If \(M=\sup_{[x_L,x_R]}g''\), the elementary interpolation estimate
\[
\operatorname{chord}_g(x_M)-g(x_M)
 \le \frac{M}{2}(x_M-x_L)(x_R-x_M)
\]
gives, because both factors on the right are at most \(1\),
\[
M>2\log\frac{67}{63}.
\]
For \(x\ge0\),
\[
\log(1+x)\ge\frac{2x}{2+x},
\]
since the difference has derivative
\(x^2/((1+x)(2+x)^2)\ge0\).  Taking \(x=4/63\) yields
\[
2\log\frac{67}{63}\ge\frac8{65}>\frac3{25}.
\]
Consequently there exists \(x_0\in(x_L,x_R)\) with
\[
g''(x_0)>\frac3{25}.
\]
By continuity there is a nondegenerate closed interval \(J\ni x_0\) such that
\[
\boxed{(\log f)''\ge\frac3{25}=2\zeta\quad\hbox{on }J},
\qquad
\zeta:=\frac3{50}.
\]

Now use Lemma~6.  On the compact interval \(J\), \(f\) has a positive minimum,
and
\[
g_h^{(j)}\longrightarrow f^{(j)}
\quad\text{uniformly},\qquad j=0,1,2.
\]
The map
\[
(u,v,w)\longmapsto \frac wu-\left(\frac vu\right)^2
\]
is continuous on \(u>0\).  Therefore there is a finite integer \(H_1\) such
that for every \(h\ge H_1\),
\[
g_h>0,\qquad
(\log g_h)''\ge\zeta=\frac3{50}
\quad\text{throughout }J.
\]

Let \(L=|J|>0\).  Since
\[
|a_h|=|a_0|\rho^h,
\]
the set \(I_h\) of integers \(k\) for which the three interpolation nodes
\(x_{h,k-1},x_{h,k},x_{h,k+1}\) all belong to \(J\) is consecutive and satisfies
\[
|I_h|\ge L|a_h|-4.
\]
Choose \(H_2\) so that \(L|a_0|\rho^h\ge8\) for \(h\ge H_2\).  Then
\[
\boxed{|I_h|\ge c_b\rho^h},
\qquad
c_b:=\frac12L|a_0|>0.
\]
For every \(k\in I_h\), with \(s_h=|a_h|^{-1}\), the exact centered
second-difference identity gives
\[
\log\frac{p_{h,k-1}p_{h,k+1}}{p_{h,k}^2}
=
\int_{-s_h}^{s_h}
 (s_h-|u|)(\log g_h)''(x_{h,k}+u)\,du
\ge
\frac{\zeta}{|a_h|^2}>0.
\]
Thus every \(k\in I_h\) is a strict failure of log-concavity.  Multiplication
by the hard-core factor \(\lambda^k/Z(T_{b,h};\lambda)\) preserves the sign of
the log-concavity determinant, so the same interval consists of failures of
the unweighted independence sequence.

Finally,
\[
\frac{m_h}{b^h}\to L_b,\qquad
\frac{\alpha_h}{b^h}\to A_b,\qquad
\frac{|a_h|}{b^h}=|a_0|q^h\to0.
\]
Since \(J\) is bounded,
\[
\sup_{k\in I_h}\left|\frac{k}{\alpha_h}-\beta_b\right|\to0,
\qquad
\beta_b=\frac{L_b}{A_b}.
\]
For \(b=64\),
\[
\beta_{64}
<
\frac{12350}{12477}
<
\frac{99}{100}.
\]
Hence there is a finite \(H_3\) such that the entire run lies below
\(0.99\,\alpha_h\) for all \(h\ge H_3\).  With
\[
H_{64}:=\max\{H_1,H_2,H_3\},
\]
all assertions of the consecutive-run theorem hold for every
\(h\ge H_{64}\).

This closes the proof without an unquantified curvature step.  The constants
\(J\) and \(H_{64}\) need not be numerically evaluated because the theorem
asserts their existence; their positivity and finiteness are established
above.
